%\input{../preamble}

%\usepackage{nomencl}

%\begin{document}

%\printnomenclature

In this section we will consider the algebraic analogues of invariant measures on groupoids, which are called \emph{invariant means}. These were originally studied in \cite{MR3448402} in the context of \emph{paradoxical decompositions}, and their connections with C*-algebras were described in \cite{MR3548134}.

We will be interested in the analytical structure that such an invariant mean will induce on a Boolean inverse monoid, and this will be revisited when we introduce sofic groupoids and study them in the next chapter. This is also a setting we will consider in the next section, on which we will prove a non-commutative version of the Loomis-Sikorski Theorem and determine Boolean semigroups which are measured semigroups of probability measure-preserving groupoids.

Although we will introduce the concepts as in \cite{MR3548134}, we will reword some definitions to fit better to the nomenclature and notation we have used so far. It should be noted that Boolean inverse monoids of \cite{MR3548134} correspond to \emph{weakly} Boolean inverse monoids as in Definition \ref{definitionbooleaninversesemigroup}.

\subsection{Covers and fixed idempotents}

\begin{definition}\label{definitioncoverinposet}\index{Cover!in a poset}
Let $L$ be a poset with zero. A \emph{cover} of an element $x\in L$ is a family $C\subseteq L$ such that whenever $e\in L$ satisfies $c\land e=0$ for all $c\in C$, then $x\land e=0$ as well. A \emph{cover} of a subset $X\subseteq L$ is a family $C\subseteq L$ which is a cover of every element of $X$.
\end{definition}

%\begin{proposition}%Using: C is a cover if whenever $e\geq c$ for all $c$ then $e\geq x$.
%If $L$ is a $\lor$-semilattice, then a finite family $C$ is a cover of $x$ if and only if $x\leq\bigvee C$.
%\end{proposition}
%\begin{proof}
%Suppose $C$ is a cover of $x$. Then $c\leq \bigvee C$ for all $c$, so $x\leq\bigvee C$. Conversely, if we assume $x\leq\bigvee C$, then whenever $p\geq c$ for all $c\in C$ yields $x\leq\bigvee C\leq p$.\qedhere
%\end{proof}



\begin{proposition}\label{propositioncoversMR3548134}
If $L$ is a generalized Boolean algebra, $X\subseteq L$ is downwards closed, and $C$ is a finite subset of $L$, then the following are equivalent:
\begin{enumerate}
\item[(1)] $C$ is a cover of $X$;
\item[(2)] For every nonzero $x\in X$ there is some $c\in C$ with $x\land c\neq 0$.
\item[(3)] $\bigvee C$ is an upper bound of $X$;
\end{enumerate}
\end{proposition}
\begin{proof}
(1)$\Rightarrow$(2): If $x\in X$ satisfies $x\land c=0$ for all $c\in C$, then the definition of cover implies $x=x\land x=0$.

(2)$\Rightarrow$(3): If $x\in X$, then the element $y=x\setminus\bigvee C$ of $X$ satisfies $y\land c=0$ for all $c\in C$, so $y=0$ by property (2) and thus $x\leq y\lor\bigvee C=\bigvee C$.

(3)$\Rightarrow$(1): If $e\in L$ satisfies $e\land c=0$ for all $c\in C$, then for all $x\in X$,
\[e\land x\leq e\land\bigvee C=\bigvee_{c\in C} e\land c=0\qedhere\]
\end{proof}

%In \cite[Definition 2.1]{MR3548134}, covers were defined as in item (2) above. Since we will only work under the hypotheses of Proposition \ref{propositioncoversMR3548134}, the definition we adopt is equivalent

We need now to consider \emph{fixed} idempotents of elements of inverse semigroups. These can be thought of as the set of fixed points of a function, as described in Example \ref{examplefixedidempotentsinix}.

\begin{definition}[{\cite[Definition 3.6]{MR3548134}}]\index{Fixed idempotent}\nomenclature[S]{$\mathcal{I}_s$}{Set of fixed idempotents of $s$}\label{definitionIsandFixs}
If $S$ is an inverse semigroup, we say that an idempotent $e\in E(S)$ is \emph{fixed} by an element $s\in S$ if $se=e$. We denote by $\mathcal{I}_s=\left\{e\in E(S):se=e\right\}$ the set of fixed idempotents of $s$.
\end{definition}

The following lemma is immediate from the facts that if $e$ is idempotent then $e^*e=e$, and that $E(S)$ is downwards directed.

\begin{lemma}\label{lemmaIs1}
If $S$ is an inverse semigroup, then for all $s\in S$,
\[\mathcal{I}_s=\left\{e\in E(S):e\leq s\right\}=\left\{e\in S:e\leq s\text{ and }e\leq s^*s\right\}.\]
In particular, $\mathcal{I}_s$ is downwards closed and $\mathcal{I}_{s^*}=\mathcal{I}_s$.
\end{lemma}

\begin{definition}[{\cite[Definition 1.7]{MR1300843}}]\nomenclature[S]{$\Fix(s)$}{Fixed point idempotent $s$}\label{definitionfixs}
If $S$ is an inverse semigroup, $s\in S$ and the meet $s\land s^*s$ exists, we denote it by $\Fix(s)=s\land s^*s$ and call $\Fix(s)$ the \emph{fixed point idempotent} of $s$.
\end{definition}

Since, by Lemma \ref{lemmaIs1}, $\mathcal{I}_s$ is the set of lower bounds of $\left\{s,s^*s\right\}$, then
\[\Fix(s)=\bigvee\mathcal{I}_s\]
in the sense that one of these elements exists if and only if the other exists, in which case they are equal and $\Fix(s)$ is the maximal idempotent fixed by $s$.

\begin{example}\label{examplefixedidempotentsinix}
If $f\in\mathcal{I}(X)$ for some set $X$, then $\Fix(f)$ is the identity function of the set of fixed points of $f$, $\left\{x\in\operatorname{dom}(f):f(x)=x\right\}$, and so under the usual interpretation that $E(\mathcal{I}(X))$ consists of subsets of $X$, $\Fix(f)$ is the set of fixed points of $f$. (See Example \ref{exampleidempotentssets}.)
\end{example}

\begin{example}
If $A\in\KB(\G)$ for some ample Hausdorff groupoid $\G$, then $\Fix(A)=A\cap\G[0]$.
\end{example}

In \cite{MR1300843}, Leech studied the \emph{fixed point operator} $\Fix$ for an inverse monoid $S$, and in particular proved that $\Fix(s)$ is defined for all $s\in S$ if and only if $S$ admits meets (with respect to its usual order). Lemma \ref{lemmaIs2}(b) below describes the same fact for inverse semigroups.

\begin{lemma}\label{lemmaIs2}
Suppose that $S$ is an inverse semigroup and $\Fix(s)$ exists for all $s\in S$. Then whenever $s,t\in S$,
\begin{enumerate}
\item[(a)] $\Fix(s)=\Fix(s^*)$;
\item[(b)] $s\land t$ exists and $s\Fix(s^*t)=t\Fix(s^*t)=s\land t$;
%\item[(c)] $s\Fix(s^*t)=s\land t$.
\item[(c)] $\Fix(s^*t)$ is a lower bound of $\left\{s^*s,t^*t\right\}$
\item[(d)] $\Fix(s^*t)=(s\land t)^*(s\land t)$;
\item[(e)] If $z\in S$ then $\Fix(s^*t)\Fix(t^*z)\leq\Fix(s^*z)$;
\item[(f)] $\Fix(st^*)=t\Fix(s^*t)t^*=s\Fix(s^*t)t^*$;
\end{enumerate}
\end{lemma}
\begin{proof}
\begin{enumerate}
\item[(a)] By \ref{lemmaIs1}, $\Fix(s)=\bigvee\mathcal{I}_s=\bigvee\mathcal{I}_{s^*}=\Fix(s^*)$.
\item[(b)] Since $s\Fix(s^*t)\leq s$ and $s\Fix(s^*t)\leq ss^*t\leq t$, then $s\Fix(s^*t)$ is a lower bound of $\left\{s,t\right\}$. If $p$ is any other lower bound of $\left\{s,t\right\}$, then $p^*p$ is an idempotent and $p^*p\leq s^*t$, so $p^*p\leq\Fix(s^*t)$, which implies $p=pp^*p\leq s\Fix(s^*t)$. This proves $s\Fix(s^*t)=s\land t$, and the other equality is similar.
%\item[(c)] By item (b), $s\Fix(s^*t)$ is a lower bound of both $s$ and $t$, so we just need to show it is the largest one: Suppose $p\leq s,t$. Then by \ref{theoremorder}(d) 
%\[(s^*t)(p^*p)=s^*(tp^*p)=s^*p=p^*p\]
%so $p^*p\in\mathcal{I}_{s^*t}$, which implies $p^*p\leq \Fix(s^*t)$, therefore
%\[p=pp^*p\leq s\Fix(s^*t)\]
%which concludes the proof that $s\Fix(s^*t)$ is the largest lower bound of $\left\{s,t\right\}$.
\item[(c)] follows from item (a) and \ref{lemmaIs1}.
\item[(d)] follows from (b) and (c).
\item[(e)] Since $\Fix(s^*t)\Fix(t^*z)\in E(S)$, then applying Lemma \ref{lemmaIs1} we obtain
\[\Fix(s^*t)\Fix(t^*z)\leq s^*tt^*z\leq s^*z\]
that is, $\Fix(s^*t)\Fix(t^*z)\in\mathcal{I}_{s^*z}$ and so $\Fix(s^*t)\Fix(t^*z)\leq\bigvee\mathcal{I}_{s^*z}=\Fix(s^*z)$
\item[(f)] By (b) and (d),
\begin{align*}
t\Fix(s^*t)t^*&=t\Fix(s^*t)\Fix(s^*t)t^*=(t\Fix(s^*t))(t\Fix(s^*t))^*\\
&=(s\land t)(s\land t)^*(s^*\land t^*)^*(s^*\land t^*)=\Fix(st^*),
\end{align*}
and the other equality is similar.\qedhere
\end{enumerate}
\end{proof}

\begin{corollary}
If $S$ is a Boolean inverse semigroup and $s,t\in S$ then $t\setminus s=t(t^*t\setminus\Fix(s^*t))$.
\end{corollary}
\begin{proof}
Using the previous lemma and Proposition \ref{propositionrelativecomplementsandoperations}(a),
\[t\setminus s=t\setminus(s\land t)=t(t^*t\setminus(s\land t)^*(s\land t))=t(t^*t\setminus\Fix(s^*t))\qedhere\]
\end{proof}

Note that, by Proposition \ref{propositioncoversMR3548134}, the condition in item (1) below is precisely condition (H) considered in \cite{MR3548134} when we restrict our study to weakly Boolean inverse semigroups (which are the semigroups of interest there).

\begin{proposition}
Given a weakly Boolean inverse semigroup $S$, the following are equivalent:
\begin{enumerate}
\item[(1)] For every $s\in S$, there is a finite subset $C\subseteq\mathcal{I}_s$ which is a cover of $\mathcal{I}_s$;
\item[(2)] For all $s\in S$, $\bigvee\mathcal{I}_s$ exists;
\item[(3)] For all $s\in S$, $\Fix(s)$ exists;
\item[(4)] $S$ is a Boolean inverse semigroup.
\end{enumerate}
\end{proposition}
\begin{proof}
(1)$\Rightarrow$(2): Suppose (1) is valid, $s\in S$ and $C\subseteq\mathcal{I}_s$ is a finite cover $C$ of $\mathcal{I}_s$. Since $\mathcal{I}_s$ is a compatible set, $\bigvee C$ exists and it is an upper bound for $\mathcal{I}_s$ by Proposition \ref{propositioncoversMR3548134}. This obviously implies $\bigvee C=\bigvee\mathcal{I}_s$.

(2)$\iff$(3) is clear from the comment below Definition \ref{definitionfixs}.

(3)$\Rightarrow$(4) follows  from Lemma \ref{lemmaIs2}(b)

(4)$\Rightarrow$ (1): If $S$ is Boolean then $\Fix(s)$ is the maximum of $\mathcal{I}_s$ and $\left\{\Fix(s)\right\}$ is a finite cover of $\mathcal{I}_s$ (by Proposition \ref{propositioncoversMR3548134}).\qedhere
\end{proof}

\subsection{Invariant means and uniform metrics}

\begin{definition}[\cite{MR3448402}]\label{definitioninvariantmean}\index{Invariant mean on a semigroup}
A \emph{mean} on a weakly Boolean inverse semigroup $S$ is a map $\mu:E(S)\to[0,\infty)$ satisfying
\begin{enumerate}
\item[(i)] $\mu(e\lor f)=\mu(e)+\mu(f)$ whenever $e,f\in E(S)$, $ef=0$.
\end{enumerate}
and a mean $\mu$ is called \emph{invariant} if
\begin{enumerate}
\item[(ii)] $\mu(s^*s)=\mu(ss^*)$ for all $s\in S$;
\end{enumerate}
A mean $\mu$ on $S$ will be called \emph{faithful} if $\mu(e)=0$ implies $e=0$, and if $S$ is a monoid and $\mu(1)=1$, we say that $\mu$ is \emph{normalized}.
\end{definition}

\begin{example}
If $X$ is a set, and again by identifying $E(\mathcal{I}(X))$ with subsets of $X$, a mean on $\mathcal{I}(X)$ corresponds to a finite, finitely additive measure on $X$ (endowed with the discrete $\sigma$-algebra).
\end{example}

\begin{example}\label{exampleinvariantmeansaremeasures}
Let $S$ be a Boolean inverse semigroup. By non-commutative Stone duality (Theorem \ref{theoremnoncommutativestoneduality}), there exists an ample Hausdorff groupoid $\G$ such that $S=\KB(\G)$. 

Identifying $E(S)$ with the collection of compact-open subsets of $\G[0]$, an invariant mean $\mu$ on $E(S)$ corresponds to a pre-measure, also denoted $\mu$, on the ring of compact-open subsets of $\G[0]$. Since $\G[0]$ is zero-dimensional, we can extend $\mu$ uniquely to a regular, locally finite Borel measure on $\G[0]$. This is a consequence of a theorem of Ma\v{r}\'\i k (\cite{MR0088532}), although a simple application of the Riesz-Markov-Kakutani Theorem also yields an elementary proof in the case at hand (see the Remark below). This interpretation is sometimes useful because we have more freedom to take unions in $\KB(\G)$, even though we do not end up with elements of $\KB(\G)$.
\end{example}

\begin{remark}
If $X$ is a zero-dimensional, locally compact Hausdorff space, then every finitely additive pre-measure $\mu$ on $\KB(X)$ (the collection of compact-open subsets of $X$) extends uniquely to a regular Borel measure\ftnt{See the beginning of Subsection \ref{subsectionl1spaces} for the definition of regular measure.} on $X$. Indeed, $\mu$ is in fact $\sigma$-additive because all sets of $\KB(X)$ are compact-open, so if $A=\bigcup_n A_n$  where $A,A_n\in\KB(X)$ and the sets $A_n$ are pairwise disjoint, then all except finitely many $A_n$ are empty, from which $\mu(A)=\sum_n\mu(A_n)$ follows.

Let $\mathcal{A}$ be the complex span of all characteristic functions of elements of $\KB(X)$, and consider the functional $\tau:\mathcal{A}\to\mathbb{C}$, $\tau(f)=\int fd\mu$.

Then $\mathcal{A}$ is dense in $C_c(X,\mathbb{C})$ (the vector space of complex-valued, compactly supported continuous functions on $X$) with uniform norm, so the continuous positive functional $\tau$ extends uniquely to a positive functional on $C_c(X,\mathbb{C})$.% The general case follows from this one, because $C_c(X)$ is a direct limit of $C_c(K)$, $K\in\KB(X)$, under canonical inclusions $C_c(K)\hookrightarrow C_c(L)$ ($K\subseteq L$), so $\tau$ still extends uniquely to a positive functional on $C_c(X)$.

By the Riesz-Markov-Kakutani Representation Theorem (\cite[Theorem 2.14]{MR924157}), $\tau$ corresponds to a unique regular Borel measure $\nu$ on $X$, which extends $\mu$ because $\nu(A)=\tau(1_A)=\mu(A)$ for all $A\in\KB(X)$.
\end{remark}

%This interpretation of invariant means as measures simplifies some calculations, since we can take arbitrary unions of elements in $\KB(\G)$, even though they are not necessarily in $\KB(\G)$. Therefore invariant measures satisfy the analog to measures.

Either using Example \ref{exampleinvariantmeansaremeasures} or proceeding in a manner similar to that of measure theory, the following properties of invariant means follow easily:

\begin{proposition}\label{propositionpropertiesofmean}
Suppose $\mu$ is a mean on a weakly Boolean semigroup $S$. Then
\begin{enumerate}
\item[(a)] $\mu(0)=0$;
\item[(b)] $\mu(e\lor f)=\mu(e)+\mu(f)-\mu(ef)$ for all $e,f\in E(S)$.
\item[(c)] If $e\leq f$ in $E(S)$ then $\mu(f)=\mu(e)+\mu(f\setminus e)$.
\end{enumerate}
\end{proposition}

Invariant means induce a metric structure on the corresponding semigroup in a natural manner.

\begin{definition}\nomenclature[S]{$d_\mu$}{Uniform metric}\index{Uniform metric}\label{definitionuniformmetric}
If $\mu$ is an invariant mean on a Boolean inverse semigroup $S$, we define the \emph{uniform metric} (induced by $\mu$) as
\[d_\mu(s,t)=\mu\left((s^*s\lor t^*t)\setminus\Fix(s^*t)\right)\]
and whenever $S$ is endowed with an invariant mean, this is the metric structure we will consider (see \ref{propositionpropertiesofuniformmetric}(b)).
\end{definition}

\begin{example}\label{uniformmetriconkbg}\nomenclature[Z]{$A\triangle B$}{Symmetric difference}
If $S=\KB(\G)$ and $\mu$ is an invariant measure on $\G[0]$, regarded as an invariant mean on $E(S)$ as in Example \ref{exampleinvariantmeansaremeasures}, then $d_\mu(A,B)=\mu(\so(A\triangle B))$, where $A\triangle B=(A\cup B)\setminus (A\cap B)$ is the symmetric difference of $A$ and $B$.
\end{example}

\begin{proposition}\label{propositionpropertiesofuniformmetric}
Suppose $\mu$ is an invariant mean on a Boolean semigroup $S$. Then if $s,t,z\in S$,
\begin{enumerate}
\item[(a)] $d_\mu(s,t)=\mu\left((s\setminus t)^*(s\setminus t)\lor(t\setminus s)^*(t\setminus s)\right)$;
\item[(b)] $d_\mu$ is a pseudometric on $S$;
\item[(c)] $d_\mu(s^{-1},t^{-1})=d_\mu(s,t)$;
\item[(d)] $d_\mu(zs,zt)\leq d_\mu(s,t)$ and $d_\mu(sz,tz)\leq d_\mu(s,t)$;
\item[(e)] $d_\mu(s_1s_2,t_1t_2)\leq d_\mu(s_1,t_2)+d_\mu(s_2,t_2)$;
\item[(f)] If $s,t,z,w\in S$ then $d_\mu(s\land t,z\land w)\leq d_\mu(s,z)+d_\mu(t,w)$;
\item[(g)] If $\left\{s,t\right\}$ and $\left\{z,w\right\}$ are compatible, then $d_\mu(s\lor t,z\lor w)\leq d_\mu(s,z)+d_\mu(t,w)$;
\item[(h)] If $s,t,z,w\in S$ then $d_\mu(s\setminus t,z\setminus w)\leq d_\mu(s,z) +d_\mu(t,w)$.
\end{enumerate}
\end{proposition}
\begin{proof}
Items (a) and (b) actually follow from Examples \ref{exampleinvariantmeansaremeasures} and \ref{uniformmetriconkbg}, but we can give a purely semigroup-theoretic proof of them:
\begin{enumerate}
\item[(a)] follows from Lemma \ref{lemmaIs2}(d) and (a);
\item[(b)] Let $s,t,z\in S$. Since $s\setminus t=\left((s\setminus t)\setminus z)\right)\lor\left((s\land z)\setminus t)\right)$, then
\[(s\setminus t)^*(s\setminus t)\leq(s\setminus z)^*(s\setminus z)\lor(z\setminus t)^*(z\setminus t)\]
and similarly with the roles of $s$ and $t$ interchanged, so simply apply (a) to obtain the result.
\item[(c)] follows from (a) and property (i) in the definition of invariant mean (\ref{definitioninvariantmean}).
\item[(d)] follows from distributivity of the product over relative complements and items (a) and (c);
\item[(e)] By (b) and (d),
\[d_\mu(s_1s_2,t_1t_2)\leq d_\mu(s_1s_2,t_1s_2)+d_\mu(t_1s_2,t_1t_2)\leq d_\mu(s_1,t_1)+d_\mu(s_2,t_2).\]
\item[(f)] Just note that
\[d_\mu(s\land t,z\land w)\leq d_\mu(s\land t, s\land w)+d_\mu(s\land w,z\land w)\leq d_\mu(t,w)+d_\mu(s,z)\]
where the second inequality follows from item (a).
\item[(g)] Using the distributive properties, we obtain
\[(s\lor t)\setminus(z\lor w)=(s\setminus(z\lor w))\lor(t\setminus(z\lor w))\leq (s\setminus z)\lor (t\setminus w)\]
and similarly $(z\lor w)\setminus(s\setminus t)\leq (z\setminus s)\lor (w\setminus t)$, then the result follows from item (a).
\item[(h)] If $S$ is a unital Boolean algebra, and if we denote the complement of an element $x\in S$ by $x^c=1\setminus x$, then
\begin{align*}
(s\setminus t)\setminus(z\setminus w)&=st^c(zw^c)^c=st^c(z^c\lor w)=st^cz^c\lor st^cw\\
&=s(t\lor w)^c\lor (sw)t^c=(s\setminus(t\lor w))\lor((sw)\setminus t)
\end{align*}
For a general Boolean inverse monoid, the same calculation on the unital Boolean algebra $\left\{x\in S:x\leq s\right\}$ (and intersecting all elements with $s$, if necessary) will also yield
\[(s\setminus t)\setminus(z\setminus w)=(s\setminus(t\lor w))\lor((s\land w)\setminus t)\]
and similarly $(z\setminus w)\setminus (s\setminus t)=(z\setminus(w\lor t))\lor((z\land t)\setminus w)$. These equalities and item (a) yield the desired inequality.\qedhere
\end{enumerate}
\end{proof}

%\end{enumerate}
%\end{proposition}
%\begin{proof}
%\begin{enumerate}
%\item[(a)] The uniform metric on $\KB(\G)$ is given by
%\[d_\mu(\alpha,\beta)=\mu(s(\alpha)\cup s(\beta))-\mu(\Fix(s^*t))=\mu(s(\alpha\cup\beta))-\mu(s(\alpha\cap\beta))=\mu(s(\alpha\cup\beta)\setminus s(\alpha\cap\beta))\]
%and since the source map is injective on both $\alpha$ and $\beta$,
%\[d_\mu(\alpha,\beta)=\mu((s(\alpha)\setminus s(\alpha\cap\beta))\cup (s(\beta)\setminus s(\alpha\cap\beta)))=\mu(s(\alpha\setminus(\alpha\cap\beta))\cup s(\beta\setminus(\alpha\cap\beta)))=\mu(s(\alpha\triangle\beta))\]
%
%The function $\mu$ is increasing, and for any three sets $A$, $B$, $C$, $A\triangle B\subseteq(A\triangle C)\cup(C\triangle B)$, so we immediately conclude that $d_\mu$ is indeed a 
%\item[(b)]
%\item[(c)]
%\end{enumerate}
%\end{proof}

%As an immediate consequence of \ref{propositionpropertiesofuniformmetric}(a) and Proposition \ref{propositionquotientrelationofbooleansemigroup}, the quotient metric space of $S$ by $d_\mu$ inherits all the structure of $S$ and $\mu$.

\begin{corollary}\label{corollaryidealofnullsetsbooleaninversesemigroup}
If $\mu$ is an invariant mean on a Boolean inverse semigroup $S$, then
\[\mathcal{N}(\mu)=\left\{n\in S:\mu(n^*n)=0\right\}=\left\{n\in S:d_\mu(0,n)=0\right\}\]
is a $\lor$-ideal of $S$, and $s\sim_{\mathcal{N}(\mu)} t$ if and only if $d_\mu(s,t)=0$. Moreover, $\mu$ factors to a faithful invariant mean on the quotient $S/\mathcal{N}(\mu)$, whose associated metric is the quotient metric of $d_\mu$.
\end{corollary}

We will still denote the invariant mean on $S/\mathcal{N}(\mu)$ by $\mu$ and the respective quotient metric by $d_\mu$.

\begin{example}
If $(\G,\mu)$ is a probability measure-preserving groupoid, and we identify $\mu$ with an invariant mean on $\BOR(\G)$ in the usual manner, then $\MEAS(\G,\mu)=\BOR(\G)/\mathcal{N}(\mu)$.
\end{example}

\begin{proposition}
Given an invariant mean $\mu$ on a Boolean inverse semigroup $S$, $d_\mu$ is a metric if and only if $\mu$ is faithful.
\end{proposition}
\begin{proof}
Suppose $d_\mu$ is a metric and $\mu(e)=0$. Then $d_\mu(e,0)=\mu(e^*e)=\mu(e)=0$, so $e=0$ and $\mu$ is faithful.

Conversely, suppose $\mu$ is faithful and $d_\mu(s,t)=0$. This means that
\[\mu(s^*s\lor t^*t\setminus\Fix(s^*t))=0,\]
so $(s^*s\lor t^*t)\setminus\Fix(s^*t)=0$. But Lemma \ref{lemmaIs2}(c) implies $\Fix(s^*t)\leq s^*s\land t^*t$, therefore $s^*s=t^*t=\Fix(s^*t)$ by faithfulness of $\mu$, thus \ref{lemmaIs2}(b) yields
\[s=ss^*s=s\Fix(s^*t)=t\Fix(s^*t)=tt^*t=t,\]
so $d_\mu$ is a metric.\qedhere
\end{proof}

The last function we will associate to an invariant mean $\mu$ on a Boolean inverse semigroup is the \emph{trace}, which is a simple extension of $\mu$ to all of $S$ and which encodes $d_\mu$ in a simpler manner.

\begin{definition}\nomenclature[S]{$\operatorname{tr}_\mu(s)$}{Trace of $s$}\index{Trace}\label{definitiontrace}
If $S$ is a Boolean inverse semigroup endowed with an invariant mean $\mu$ and $s\in S$, define the \emph{trace} of $s$ as $\tr(s)=\mu(\Fix(s))$ (or $\tr_\mu$ if we need to make $\mu$ explicit).
\end{definition}

\begin{example}
Let $(X,\mu)$ be a standard probability space and $S$ the inverse semigroup of measure-preserving Borel automorphisms of $X$. By identifying $E(S)$ with the set of Borel subsets of $X$ and $\mu$ with an invariant mean on $S$, the trace of $f\in S$ is
\[\tr_\mu(f)=\mu\left(\left\{x\in X:f(x)=x\right\}\right).\]
\end{example}

\begin{example}
As a particular case of the previous example, let $X$ be a finite set with $n$ elements (seen as a discrete Borel space) and $\mu_\#$ the normalized counting measure on $X$: $\mu_\#(A)=\#A/n$ for all $A\subseteq X$, where $\#A$ stands for the cardinality of $A$. Then $\mathcal{I}(X)$ is precisely the set of measure-preserving Borel automorphisms of $(X,\mu_\#)$.

Let $S$ be the inverse semigroup of $n\times n$ partial permutation matrices, that is, $n\times n$ matrices such that every row and every column has entries in $\left\{0,1\right\}$, and at most one entry in each row and in each column is nonzero. Then $S$ and $\mathcal{I}(X)$ are isomorphic via the map $\theta:\mathcal{I}(X)\to S$, $f\mapsto\theta(f)=[f_{ij}]_{ij}$ given by
\[f_{ij}=\begin{cases}
1,&\text{if }j\in\operatorname{dom}(f)\text{ and }f(j)=i,\\
0,&\text{otherwise.}
\end{cases}\]

Composing the trace on $\mathcal{I}(X)$ (induced by $\mu_\#$) with $\theta^{-1}$, we obtain the usual normalized matrix trace on $S$: $\tr([m_{ij}]_{ij})=(1/n)\sum_{i=1}^n m_{ii}$.
\end{example}

\begin{example}
If $S=\KB(\G)$ and $\mu$ is an invariant mean on $S$, then $\tr(A)=\mu(A\cap\G[0])$ for all $A\in S$.
\end{example}

\begin{theorem}\label{theoremtracepreservingisisometric}
Let $S,T$ be Boolean inverse monoids with faithful and normalized invariant means $\mu$ and $\nu$, respectively. Then a semigroup morphism $\theta:S\to T$ preserves traces if and only if it is isometric with respect to $d_\mu$ and $d_\nu$. In this case, $\theta$ is a morphism of Boolean inverse monoids satisfying $\theta(1)=1$ (i.e., $\theta$ is proper).
\end{theorem}
\begin{proof}
First suppose $\theta$ preserves the trace. Then the definitions of $d_\mu$ and $\tr_\mu$, and Proposition \ref{propositionpropertiesofmean}(c) imply
\[d_\mu(s,t)=\tr_\mu(s^*s)+\tr_\mu(t^*t)-\tr_\mu(s^*st^*t)-\tr_\mu(s^*t)\]
and similarly for $d_\nu$, and so it is clear that $\theta$ is isometric.

Conversely, suppose $\theta$ is isometric. First we prove that $\theta(0)=0$ and $\theta(1)=1$. Since $0\leq 1$ then $\theta(0)\leq \theta(1)$, so
\[1=d_\mu(0,1)=d_\nu(\theta(0),\theta(1))=\nu((\theta(0)\lor\theta(1))\setminus(\theta(0)\theta(1)))=\nu(\theta(1)\setminus\theta(0))\]
and faithfulness of $\nu$ implies $\theta(1)\setminus\theta(0)=1$, so $\theta(1)=1$ and $\theta(0)=0$. Now note that for all $s\in S$,
\[\tr_\mu(s)=\mu(\Fix(s))=\mu(s^*s)-\mu\left(s^*s\setminus\Fix(s)\right)=d_\mu(s^*s,0)-d_\mu(s^*s,s)\]
and similarly for $\tr_\nu$. Therefore $\theta$ preserves traces.

In this case, we have already verified that $\theta(0)=0$ and $\theta(1)=1$. In order to verify that $\theta$ preserves meets, we first verify that $\theta(\Fix(s))=\Fix(\theta(s))$ for all $s\in S$: Indeed, $\Fix(s)$ is an idempotent fixed by $s$, so $\theta(\Fix(s))$ is an idempotent fixed by $\theta(s)$ and so $\theta(\Fix(s))\leq\bigvee\mathcal{I}_{\theta(s)}=\Fix(\theta(s))$ (from Definition \ref{definitionfixs} and the paragraph following it).  As
\begin{align*}
\nu(\theta(\Fix(s)))&=\tr_\nu(\theta(\Fix(s)))=\tr_\mu(\Fix(s))=\mu(\Fix(s))=\tr_\mu(s)=\tr_\nu(\theta(s))\\
&=\nu(\Fix(\theta(s))),
\end{align*}
faithfulness of $\nu$ implies $\theta(\Fix(s))=\Fix(\theta(s))$.

Given $s,t\in S$, since $s\land t=s\Fix(s^*t)$ (Lemma \ref{lemmaIs2}(b)), it follows from the previous argument that $\theta(s\land t)=\theta(s)\land \theta(t)$.

Similarly, in order to prove that $\theta$ is a $\lor$-morphism we first assume $e,f\in E(S)$. Then $\theta(e\lor f)$ is an upper bound of $\left\{\theta(e),\theta(f)\right\}$, and so $\theta(e)\lor\theta(f)\leq\theta(e\lor f)$. But on the other hand,
\begin{align*}
\nu(\theta(e)\lor\theta(f))&=\tr_\nu(\theta(e))+\tr_\nu(\theta(f))-\tr_\nu(\theta(ef))\\
&=\tr_\mu(e)+\tr_\mu(f)-\tr_\mu(ef)\\
&=\tr_\mu(e\lor f)=\tr_\nu(\theta(e\lor f))=\nu(\theta(e\lor f))
\end{align*}
and again faithfulness of $\nu$ implies $\theta(e)\lor\theta(f)=\theta(ef)$.

Given compatible elements $s,t\in S$, $\theta(s)$ and $\theta(t)$ are compatible in $T$, and $\theta(s)\lor\theta(t)\leq\theta(s\lor t)$. But
\begin{align*}
(\theta(s)\lor\theta(t))^*(\theta(s)\lor\theta(t))&=\theta(s)^*\theta(s)\lor\theta(t)^*\theta(t)=\theta(s^*s)\lor\theta(t^*t)=\theta(s^*s\lor t^*t)\\
&=\theta((s\lor t)^*(s\lor t))=\theta(s\lor t)^*\theta(s\lor t)
\end{align*}
so $\theta(s)\lor\theta(t)=\theta(s\lor t)$ (Theorem \ref{theoremorder}(e)).\qedhere
\end{proof}
%\begin{proposition}\label{propositionmetric}
%Given $,\beta,\gamma,\delta\in [[G]]_B$;
%\begin{enumerate}
%\item $d_{\mu}(\alpha,\beta)=d_{\mu}(\alpha^{-1},\beta^{-1})$.
%%%%%%%\item $d_\mu(\alpha,\beta)=\mu(s(\alpha)\triangle s(\beta))+\mu\left\{x\in s(\alpha)\cap s(\beta):s|_\alpha^{-1}(x)\neq s|_\beta^{-1}(x)\right\}$;
%\item $d_{\mu}(\alpha\beta,\gamma\delta)\leq d_{\mu}(\alpha,\gamma)+d_\mu(\beta,\delta)$;
%\item $d_{\mu}(\alpha,\beta^{-1})\leq d_{\mu}(\alpha,\alpha\beta\alpha)+d_{\mu}(\beta,\beta\alpha\beta)$;
%\end{enumerate}
%\end{proposition}
%\begin{proof}
%\begin{enumerate}
%\item Since $\mu(s(\alpha\setminus\beta))=\mu(r(\alpha^{-1}\setminus\beta^{-1}))=\mu(s(\alpha^{-1}\setminus\beta^{-1})$ and similarly with the roles of $\alpha$ and $\beta$ replaced. Add these equalities and we are done.\qedhere
%\item The sets on the right-hand side for which we take their measures form a decomposition of $s(\alpha\triangle\beta)$.
%\item PROOF 1: We have $AB\setminus CD=(A\setminus C)B\cup(A\cap C)(B\setminus D)$. Since $s((A\cap C)(B\setminus D))\subseteq s(B\setminus D)$ and $r((A\setminus C)B)\subseteq r(A\setminus C)$, we obtain
%\begin{align*}
%\mu(s(AB\setminus CD))&\leq\mu(r(A\setminus C))+\mu(s(B\setminus D))=\mu(s(A\setminus C))+\mu(s(B\setminus D))
%\end{align*}
%Similarly, $\mu(s(CD\setminus AB))\leq\mu(s(C\setminus A))+\mu(s(D\setminus B))$, so adding there two inequalities yields the result.%PROOF 2:
%\item Note that $\alpha\beta\setminus\alpha\delta=\alpha(\beta\setminus\delta)$ for all $\alpha,\beta\delta\in[[G]]_B$. From this we can infer, taking sources, that $d_\mu(\alpha\beta,\alpha\delta)\leq d_\mu(\beta,\delta)$. Similarly, $d_\mu(\alpha\delta,\gamma\delta)\leq d_\mu(\alpha,\gamma)$, so
%\[d_\mu(\alpha\beta,\gamma\delta)\leq d_\mu(\alpha\beta,\alpha\delta)+d_\mu(\alpha\delta,\gamma\delta)\leq d_\mu(\beta\delta)+d_\mu(\alpha,\gamma).\]
%\item One first verifies that
%\[s(\alpha\triangle\alpha\beta\alpha)=s(\alpha)\setminus r(\beta)\cup\left\{x\in s(\alpha)\cap r(\beta):s|_\alpha^{-1}x\neq\left(r|_\beta^{-1}x\right)^{-1}\right\}\]
%and the two sets on the right-hand side are disjoint. Taking their measures yields
%\[\mu(s(\alpha)\setminus r(\beta))+\left\{x\in s(\alpha)\cap r(\beta):s|_\alpha^{-1}x\neq\left(r|_\beta^{-1}(x)\right)^{-1}\right\}=d_\mu(\alpha,\alpha\beta\alpha)\]
%and an application of the second item finishes the proof.\qedhere
%\end{enumerate}
%\end{proof}

%\end{document}